\chapter{API Quick Reference: Ecosystem and Platforms}
\label{ch:api-quick-reference-vol2}

Appendix~\ref{ch:api-quick-reference} did the same job for the core framework and graphics
classes; this appendix extends that lookup companion to the namespaces covered in
Parts~V--IX --- Input, Audio, Media, Devices and Sensors, Networking, GamerServices, and a
handful of \repolink{sharp-runtime} core types every other chapter in those parts leans on. As in
Appendix~\ref{ch:api-quick-reference}, every entry here was read directly from the
corresponding header at the time of writing, and this is a summary of classes already given
full narrative treatment elsewhere in this book, not a substitute for reading that treatment.

\section{Input: Keyboard, Mouse, GamePad, and Touch}
\label{sec:api-input}

\cnaclass{Keyboard} and \cnaclass{Mouse} are both non-instantiable static classes exposing a
single \texttt{GetState()} that returns an immutable snapshot --- \cnaclass{KeyboardState} (an
\texttt{IsKeyDown} / \texttt{IsKeyUp} / \texttt{GetPressedKeys()} query surface plus an indexer
returning a \cnaclass{KeyState}) and \cnaclass{MouseState} (\texttt{X} / \texttt{Y},
\texttt{LeftButton} / \texttt{MiddleButton} / \texttt{RightButton} / \texttt{XButton1} / \texttt{XButton2},
and \texttt{ScrollWheelValue}) respectively. Both classes carry a cluster of \texttt{CNAEXT/EXT}
extensions that fill gaps a desktop port genuinely needs: \cnaclass{Keyboard}'s
\texttt{GetModStateEXT} / \texttt{GetScancodeNameEXT} / \texttt{GetKeyNameEXT} (and their inverses)
expose layout-independent versus layout-dependent key naming that XNA's console-era API never
needed to distinguish; \cnaclass{Mouse}'s \texttt{SetCaptureEXT} / \texttt{GetGlobalPositionEXT}/%
\texttt{WarpGlobalEXT} and relative-mouse-mode property support desktop window-manager
interactions with no Xbox~360 equivalent at all. \cnaclass{MouseState} additionally carries a
\texttt{CNAEXT GetHorizontalScrollWheelValueEXT} property --- deliberately excluded from
\texttt{Equals} / \texttt{GetHashCode} so those two stay byte-identical to FNA's own, vertical-only
comparison, even though the field itself is real and populated from SDL's horizontal wheel axis.

\textbf{Coordinate warning.} Absolute mouse/button/touch reads and
\cnaclass{Mouse}::\texttt{SetPosition} do not have one uniform logical-presentation contract.
SDL-renderer products have the complete mapping; EasyGL/Canvas cover only their default
height-driven mode; WebGPU/SDL GPU mix window and physical-pixel spaces on high-density
displays; Vulkan passes through after a non-1:1 resize; and FREEDIRECT changes behavior after its
first \texttt{Present()}. See \S\ref{sec:input-logical-coordinate-routing} and the complete
renderer matrix in \S\ref{sec:backend-input-coordinate-matrix} before treating
\texttt{Mouse::GetState()} coordinates as rendered logical pixels.

\cnaclass{GamePad} follows the same static-class shape, with \texttt{GetCapabilities}/%
\texttt{GetState} (two overloads, the second taking an explicit \cnaclass{GamePadDeadZone} mode)
and \texttt{SetVibration}. Its \cnaclass{CNAEXT} surface is the largest of any class in this
appendix, and reflects a real hardware gap rather than a naming inconsistency: XNA's Xbox~360
gamepad model has no vocabulary at all for a modern controller's light bar, trigger vibration
motors, touchpad, gyroscope, accelerometer, back paddles, or the several controller identity
queries (\texttt{GetGUIDEXT}, \texttt{GetSerialEXT}, \texttt{GetSteamHandleEXT}, and more) that a
2020s USB/Bluetooth controller stack makes available and an Xbox~360 pad simply didn't have.
\cnaclass{GamePadCapabilities} mirrors this exactly: roughly twenty real-XNA
\texttt{HasXxxButton} / \texttt{HasXxxThumbStick} / \texttt{HasXxxTrigger} flags, plus ten further
\texttt{CNAEXT}-tagged \texttt{HasXxxEXT} flags for exactly the modern features above.

\cnaclass{Touch::TouchPanel} is, like \cnaclass{Keyboard} / \cnaclass{Mouse} / \cnaclass{GamePad}, a
static class --- \texttt{GetState()} returns a \cnaclass{Touch::TouchCollection} (an
STL-container-shaped, technically-always-mutable collection that nonetheless reports
\texttt{IsReadOnly} as \texttt{true}, faithfully matching FNA's own documented inconsistency
rather than \\ "fixing" it) and \texttt{ReadGesture()} dequeues the next
\cnaclass{Touch::GestureSample}. \texttt{DisplayWidth} / \texttt{DisplayHeight}/%
\texttt{DisplayOrientation} and \texttt{EnabledGestures} round out the public surface; a cluster
of \texttt{CNAEXT} methods (\texttt{SetFinger}, \texttt{INTERNAL\_onTouchEvent}, \texttt{Update})
mark the platform-input-bridge plumbing FNA itself keeps \texttt{internal}, exposed here only
because C\texttt{++} has no assembly-internal visibility modifier to fall back on.

\section{Audio and Media: SoundEffect, Song, and MediaPlayer}
\label{sec:api-audio-media}

\cnaclass{SoundEffect} is move-only (matching FNA's own single-object instance-tracking
contract, where disposing the source effect cascades to every \cnaclass{SoundEffectInstance} it
created) and offers three constructor families: loading a file directly
(\texttt{CNAEXT}, since real XNA never constructs a \cnaclass{SoundEffect} from a bare path
outside the content pipeline), and two raw-PCM-buffer constructors --- one plain, one with
explicit loop points --- that both require genuinely headerless, signed 16-bit PCM samples
rather than a whole WAV file's bytes, a documented and easy mistake to make. \texttt{Play()}
(two overloads: fire-and-forget, or with explicit volume/pitch/pan) and \texttt{CreateInstance()}
are the two ways to actually hear it; four static properties (\texttt{MasterVolume},
\texttt{DistanceScale}, \texttt{DopplerScale}, \texttt{SpeedOfSound}) govern process-wide 3D-audio
approximation, since SDL3\_mixer itself has no native 3D audio model to fall back on.
\cnaclass{SoundEffectInstance} is the mutable playback handle \texttt{CreateInstance()} returns:
\texttt{Play} / \texttt{Stop} / \texttt{Pause} / \texttt{Resume}, \texttt{Volume} / \texttt{Pan}/%
\texttt{Pitch} / \texttt{IsLooped} properties, and \texttt{Apply3D} (an attenuation-and-pan
approximation over listener/emitter positions, not a real 3D mixer feature). Its subclass
\cnaclass{DynamicSoundEffectInstance} adds a \texttt{BufferNeeded} event and
\texttt{SubmitBuffer} (plus a \texttt{CNAEXT SubmitFloatBufferEXT} float32 variant) for
streaming/procedural audio that has no pre-loaded \cnaclass{SoundEffect} behind it at all.

\cnaclass{Song} is a content handle rather than a playback controller --- \texttt{Name},
\texttt{Duration}, \texttt{Album} / \texttt{Artist} / \texttt{Genre} (all \texttt{nullptr} unless the
song came from a real \cnaclass{MediaLibrary} scan rather than a bare file path), and real,
tag-derived \texttt{IsRated} / \texttt{Rating} / \texttt{TrackNumber} properties populated from a
file's actual ID3v2/Vorbis metadata rather than left as permanent stubs. \cnaclass{MediaPlayer}
is, like the input classes above, a static class: \texttt{Play} (three overloads: single song,
a \cnaclass{SongCollection}, or a collection plus a starting index), \texttt{Pause}/%
\texttt{Resume} / \texttt{Stop} / \texttt{MoveNext} / \texttt{MovePrevious}, and
\texttt{IsRepeating} / \texttt{IsShuffled} / \texttt{IsMuted} / \texttt{Volume} / \texttt{PlayPosition}
properties, plus an \texttt{ActiveSongChanged} and \texttt{MediaStateChanged} event pair.

\section{Devices and Sensors: Accelerometer and Compass}
\label{sec:api-devices-sensors}

Chapter~\ref{ch:devices-sensors} gives this namespace's full narrative treatment, including its
concurrency findings; this section is the lookup-companion summary. \cnaclass{Accelerometer} and
\cnaclass{Compass} both derive from the template base \cnans{SensorBase<TSensorReading>}
(\cnans{TSensorReading} is \cnaclass{AccelerometerReading} or \cnaclass{CompassReading}
respectively), which supplies every sensor class's shared surface: \texttt{CurrentValue} (throws
\texttt{System::InvalidOperationException} unless \texttt{IsSupported} is \texttt{true} ---
callers must check \texttt{IsSupported} first, never catch the exception as a substitute),
\texttt{IsDataValid}, \texttt{TimeBetweenUpdates} (a \cnaclass{System::TimeSpan}, real WP7 default
2~milliseconds, re-read fresh on every incoming sample so a change while the sensor is running
takes effect immediately without a \texttt{Stop()} / \texttt{Start()} cycle), the
\texttt{CurrentValueChanged} event, and the pure-virtual \texttt{Start()} / \texttt{Stop()} pair
every derived class implements against its own hardware.

\cnaclass{Accelerometer} adds a static \texttt{IsSupported} and a \texttt{State} property ---
the one strict-XNA-API member of \cnaclass{SensorState} among this namespace's four sensor
classes, per that enum's own documentation; \cnaclass{Compass} / \cnaclass{Gyroscope} / \texttt{Motion}
expose the identical property only as a \texttt{CNAEXT} symmetry addition, since the real WP7
\cnaclass{Compass} class has no \texttt{State} member at all. \cnaclass{Accelerometer} also
carries a legacy \texttt{ReadingChanged} event, deprecated on real WP7 7.1 in favor of
\texttt{CurrentValueChanged} but still present and still raised here for API completeness with
every real WP7 version; \texttt{CurrentValueChanged} always fires first and
\texttt{ReadingChanged} second for the same reading, never the reverse and never interleaved.
Each delivered \cnaclass{AccelerometerReading} carries a \texttt{Timestamp}
(\cnaclass{System::DateTimeOffset}) and an \texttt{Acceleration} (\cnaclass{Vector3}, in
$\text{m/s}^2$, SDL's own native unit for this axis).

\cnaclass{Compass} has no SDL-backed implementation on any platform --- SDL3 exposes no
magnetometer API anywhere --- so its real data path exists only on Android, via a dedicated
native backend (\cnans{Detail::AndroidCompassBackend}) reading the NDK's rotation-vector and
magnetic-field sensors directly, with no SDL involvement at all. On every other supported
platform, \texttt{IsSupported} is permanently \texttt{false} and \texttt{Start()} always throws
\cnaclass{SensorFailedException}. Each delivered \cnaclass{CompassReading} carries
\texttt{HeadingAccuracy} / \texttt{MagneticHeading} / \texttt{TrueHeading} (all
\texttt{double}, in degrees) and \texttt{MagnetometerReading} (\cnaclass{Vector3}, in
micro-teslas). \texttt{TrueHeading} is honestly documented as currently identical to
\texttt{MagneticHeading} on every backend in this codebase: a genuine true-north correction needs
magnetic declination, which in turn needs geographic location data this namespace does not yet
implement, and this project's own methodology forbids fabricating an assumed declination rather
than reporting the most accurate value actually available. \cnaclass{Compass} additionally
raises a \texttt{Calibrate} event (\cnaclass{CalibrationEventArgs}) when the backend detects the
sensor needs recalibration.

\cnaclass{SensorState} (\texttt{NotSupported} / \texttt{Ready} / \texttt{Initializing} /
\texttt{NoData} / \texttt{NoPermissions} / \texttt{Disabled}) is shared across all four sensor
classes in this namespace, but only four of its six values are ever actually produced by any of
them today, confirmed by reading every assignment site directly: \texttt{NoData} and
\texttt{NoPermissions} are declared and reachable in principle but currently dead code across
this entire codebase, documented as such rather than silently left unexplained.

\section{Networking and GamerServices}
\label{sec:api-net-gamerservices}

\cnaclass{NetworkSession} is the largest single class in this appendix, and its size is
representative of what it actually does: five static factory families
(\texttt{Create} / \texttt{Find} / \texttt{Join} / \texttt{JoinInvited}, each with a synchronous form
and a \texttt{Begin} / \texttt{End} fake-async pair, per the pattern Chapter~\ref{ch:storage} introduces for
\cnaclass{Storage}), four gamer collections (\texttt{AllGamers}/%
\texttt{LocalGamers} / \texttt{RemoteGamers} / \texttt{PreviousGamers}), and six events
(\texttt{GamerJoined} / \texttt{GamerLeft} / \texttt{HostChanged} / \texttt{GameStarted}/%
\texttt{GameEnded} / \texttt{SessionEnded}). \texttt{AllowHostMigration} and
\texttt{SimulatedLatency} / \texttt{SimulatedPacketLoss} are worth flagging precisely because they
invert this book's usual ``faithful shape, inert behavior'' pattern: FNA declares all three as
plain, stored-but-unused properties with zero backing logic, while CNA's own \texttt{ENetBackend}
gives every one of them a real implementation --- genuine deterministic host migration, and
genuine probabilistic latency/loss injection scoped to application data only. \cnaclass{NetworkGamer}
(\texttt{Id}, \texttt{IsHost}, \texttt{IsLocal}, \texttt{IsReady}, \texttt{RoundtripTime}) is a
second, smaller example of the same pattern: several of its properties (\texttt{Id} always~0,
\texttt{IsHost} always~\texttt{true}) are permanently hardcoded stubs in FNA itself, restored to
real, wire-negotiated values here via \texttt{CNAEXT} setters FNA has no equivalent for.

\cnaclass{LeaderboardReader} provides paged access to a leaderboard via three static
\texttt{Read} overloads (a plain page, a page centered on a pivot gamer, or a page restricted to
an explicit gamer set) each with a synchronous and \texttt{Begin} / \texttt{End} form,
\texttt{PageUp} / \texttt{PageDown}, and \texttt{CanPageUp} / \texttt{CanPageDown} / \texttt{Entries}
properties --- real, local-store-backed behavior sorted by rating descending, where FNA's own
equivalent is uniformly a \texttt{NotSupportedException} stub. \cnaclass{Achievement} is a small,
immutable value type (\texttt{Key} / \texttt{Name} / \texttt{Description} / \texttt{GamerScore}/%
\texttt{IsEarned} / \texttt{EarnedDateTime}); its one XNA method, \texttt{GetPicture()}, is a
deliberate, documented exception to this book's ``restore what FNA stubbed out'' pattern ---
Xbox~360 achievement artwork was streamed from Xbox LIVE at request time with no real local
equivalent to substitute, so it stays throwing by design, the same reasoning Chapter~\ref{ch:gamerservices}'s
\texttt{Guide} coverage already applies to its own Xbox-Live-only methods.

\section{Storage: StorageDevice and StorageContainer}
\label{sec:api-storage}

\cnaclass{StorageDevice} exposes \texttt{FreeSpace} / \texttt{TotalSpace} (both real values read
from the underlying filesystem) and \texttt{IsConnected}, following the XNA~4.0 fake-async
pattern throughout: \texttt{BeginOpenContainer} / \texttt{EndOpenContainer} open a container by
display name; \texttt{DeleteContainer} removes one outright; and a four-overload
\texttt{BeginShowSelector} / \texttt{EndShowSelector} pair (with or without an explicit
\cnaclass{PlayerIndex}, with or without a required free-space/directory-count) surfaces the
device-selection UI real XNA needed on a physical Xbox~360, even though this implementation
only ever resolves to the local filesystem. Two \texttt{CNAEXT} static methods,
\texttt{SetAppNameEXT} / \texttt{GetStorageRootEXT}, are the real root every other call in this
namespace ultimately resolves against --- a game is expected to call \texttt{SetAppNameEXT}
once, early, before any other storage call. \cnaclass{StorageContainer} provides the actual
file/directory operations relative to that root: \texttt{CreateDirectory}/%
\texttt{DirectoryExists} / \texttt{DeleteDirectory}, \texttt{GetDirectoryNames} (with a
glob-pattern overload), \texttt{CreateFile} / \texttt{FileExists} / \texttt{DeleteFile},
\texttt{GetFileNames} (same glob-pattern overload), and a three-overload \texttt{OpenFile}
family (mode alone; mode plus access; mode plus access plus sharing), each returning a real
\cnaclass{System::IO::Stream}. Every relative path passed to any of these methods is resolved
by simple path concatenation with no \texttt{..}-traversal guard of any kind --- Chapter~\ref{ch:storage}
covers this real, currently open gap in full; a filename that ultimately came from outside a
game's own hardcoded naming scheme should be validated by the caller before it ever reaches
this API.

\section{SharpRuntime Core Types Used Everywhere}
\label{sec:api-sharpruntime-core}

Three \repolink{sharp-runtime} types appear in nearly every class this book names, and are worth
a quick-reference entry of their own rather than re-explaining at every use site.
\cnaclass{System::TimeSpan} stores a tick count (100~nanoseconds each, matching .NET exactly)
behind \texttt{Days} / \texttt{Hours} / \texttt{Minutes} / \texttt{Seconds} / \texttt{Milliseconds} and
their \texttt{TotalXxx} fractional counterparts, \texttt{FromXxx} static factories for every
unit, and a full arithmetic/comparison operator set --- the type Chapter~\ref{ch:game-loop}'s
\cnaclass{GameTime} and the \cnaclass{Storage}, \cnaclass{NetworkSession}, and
\cnaclass{SoundEffect} chapters all build on for durations and timestamps.
\cnaclass{System::EventHandler<TEventArgs>} is a deliberate architectural departure from real
\uppercase{.}NET worth stating plainly: in C\#, \texttt{EventHandler<T>} is only a delegate
\emph{type}, with the \texttt{event} keyword itself supplying subscription storage and
invocation; C\texttt{++} has no \texttt{event} keyword, so this class folds both roles into one
--- \texttt{operator+=} / \texttt{Add} (returning a removal \texttt{Token}) to subscribe,
\texttt{Remove} / \texttt{Clear} to unsubscribe, and \texttt{Raise} / \texttt{Invoke} to fire every
subscriber against a snapshot of the handler list (so a handler that unsubscribes itself
mid-raise cannot corrupt the in-progress dispatch). A \texttt{CNAEXT SetReplayHook} exists
specifically to model the handful of real XNA events --- \cnaclass{NetworkSession::GamerJoined}
chief among them --- that are documented to replay already-happened state to a new subscriber
immediately upon \texttt{+=}, a quirk \texttt{EventHandler<T>} cannot reproduce automatically on
its own. \cnaclass{System::IO::Stream} is a lightweight subset of .NET's own \texttt{Stream}:
pure-virtual \texttt{Read} / \texttt{Close} / \texttt{Length}, with \texttt{Write} / \texttt{WriteByte}/%
\texttt{Seek} / \texttt{SetLength} defaulting to \texttt{NotSupportedException} for a read-only
stream and available to override where writing or seeking is genuinely supported --- the base
class every concrete stream this book names (\cnaclass{MemoryStream}, and the real filesystem
streams \cnaclass{StorageContainer::OpenFile} in Chapter~\ref{ch:storage} returns) derives from.

\cnaclass{MemoryStream} is the concrete in-memory \cnaclass{Stream} this ecosystem uses when
there is no meaningful file path (notably loaded Android content). Its default constructor makes
an empty writable byte vector. The raw-buffer constructor deliberately differs from .NET:
it defaults to read-only, where real \texttt{MemoryStream(byte[])} defaults to writable. Calling
\texttt{Close()} rejects subsequent normal stream operations (including \texttt{Read},
\texttt{Write}, \texttt{Seek}, \texttt{Length}, and \texttt{Position}), but intentionally leaves
\texttt{ToArray()} and \texttt{GetBuffer()} available for post-close extraction. The former
returns an independent vector; the latter exposes the live backing vector and so cannot be held
across a write that may reallocate it. Argument validation rejects null buffers and negative
offsets/counts, while a write whose \texttt{Position + count} overflows the runtime's signed
integer range throws \cnaclass{System::IO::IOException} rather than corrupting memory --- Ch.\ref{ch:sharp-runtime-namespaces}.

\section{Reading this appendix safely}

The same caution Appendix~\ref{ch:api-quick-reference} states applies here without
modification: every
class name above is exactly the name to search for in the corresponding repository's own
\texttt{include/} tree, and this table is a snapshot of one point in the project's history, not
a live mirror of it.
